%% table 객체: tab:prior-vs-thesis — 교전 단위 예측의 기존 정의(Baek & Kwon 2026)와 본 논문의 처리 (근거: MASTER_THESIS_RESEARCH_PLAN §3 표를 3인칭·처리 중심으로 재구성)
\begin{table}[!t]
  \centering
  \caption{교전 단위 예측의 기존 정의와 본 논문의 처리}
  \label{tab:prior-vs-thesis}
  \footnotesize
  \begin{tabularx}{\linewidth}{@{}>{\raggedright\arraybackslash}p{1.7cm}>{\raggedright\arraybackslash}p{3.6cm}>{\raggedright\arraybackslash}p{3.3cm}X@{}}
    \toprule
    구성 요소 & 기존 정의 \cite{baek2026killconditioned} & 남은 문제 & 본 논문의 처리 (장) \\
    \midrule
    교전 단위 & 게임의 이유로 정한 시간·공간 상수의 킬 군집; 첫 킬 \cogLead{}초 전 존재 게이트 & 상수의 경험적 근거와 민감도 & 코퍼스에서 추정한 $G$, $D$; 규칙에 고정한 $R$, $B$; 상수별 출처 분류; 사례 구성 민감도 (제 \ref{ch:engagement} 장) \\
    결과 가치 & 킬·오브젝트·현상금에 \cogWeights{}개 가중을 곱한 교환가치의 부호 (교전 국소) & 경기의 최종 승패와의 연결 & 동결 승률 평가기의 교전 구간 변화 방향 $\SVI$; 측정의 검증 (제 \ref{ch:value} 장) \\
    사전 예측 & \cogWindow{}초 관측 창; 표현$\times$학습기 조합 비교 & 현재 승률·시간을 넘는 이득의 기준선; 입력의 사후 정보 배제 & 고정 사양 $q$ 대 $\PTflex$의 동일 행 짝 대비; 정보군 비교; 예측 시점 이전 정보만의 입력 규칙 (제 \ref{sec:data-inputs} 절, 제 \ref{ch:prediction} 장) \\
    검증 & 초기 골드·시간 층화; 인접 패치 평가 & 균형 상태, 작은 변화, 확률 품질, 관측 갱신, 다른 지역에서의 범위 & $\Bforty$, 작은 변화 제외, 점수 분해, 관측 갱신 층화, 점수 전용 외부 평가와 부트스트랩 (제 \ref{ch:validation} 장) \\
    \bottomrule
  \end{tabularx}
\end{table}
